\documentclass[10pt,a4paper]{amsart}
\usepackage[margin=19mm]{geometry}
\usepackage{amsmath,amssymb,mathtools}
\usepackage[scr=rsfs,cal=euler]{mathalfa}
\usepackage{xfrac}
\usepackage[unicode,hidelinks,bookmarks=false]{hyperref}
\newcommand{\ie}{i.e.\ }
\newcommand{\cf}{cf.\ }
\newcommand{\resp}{respectively\ }
\newcommand{\Subsection}{\subsection}
\providecommand{\nobreakdash}{\nobreak\hbox{-}}
\setlength{\emergencystretch}{2em}
\multlinegap0pt
\usepackage{amssymb}
\usepackage{siunitx}
\usepackage{mathtools}
\usepackage[shortlabels]{enumitem}
\usepackage{dsfont}
\usepackage{bm}
\newtheorem{assumption}{Assumption}
\newtheorem{proposition}{Proposition}
\newcommand{\R}{\mathbb{R}}
\newcommand{\antoine}[1]{{\textcolor[rgb]{1,0,0}{\textsf{[Antoine: {#1}]}}}}
\newcommand{\constraintSet}{{S_{\mathcal{C}}}}
\newcommand{\diam}{\text{diam}}
\newcommand{\dynamicsFullState}{f}
\newcommand{\dynamicsPGD}{f_{\text{PGD}}}
\newcommand{\frank}[1]{{\textcolor[rgb]{0,0,1}{\textsf{[Frank: {#1}]}}}}
\newcommand{\hessianBound}{\mu}
\newcommand{\iterationHorizon}{N}
\newcommand{\linearizationError}{r}
\newcommand{\obj}{J}
\newcommand{\outputs}{y}
\newcommand{\param}{\theta}
\newcommand{\paramDim}{d}
\newcommand{\paramNominal}{\hat\param}
\newcommand{\paramSet}{\Theta}
\newcommand{\proj}{\text{proj}}
\newcommand{\ra}{\rightarrow}
\newcommand{\ratePGD}{\alpha}
\newcommand{\ratePGDError}{\Delta\ratePGD}
\newcommand{\ratePGDMax}{C_\ratePGD}
\newcommand{\ratePGDMin}{c_\ratePGD}
\newcommand{\ratePGDNominal}{\hat\ratePGD}
\newcommand{\smoothnessConstant}{L}
\newcommand{\state}{\zeta}
\newcommand{\stateError}{\Delta\state}
\newcommand{\stateNominal}{\hat\state}
\newcommand{\strongConvexityConstant}{m}
\newcommand{\var}{\xi}
\newcommand{\varConvexHull}{S_{\var, \text{CH}}}
\newcommand{\varDim}{n}
\newcommand{\varForwardReachableSet}[1]{S_{\var, {#1}}}
\newcommand{\varInitialSet}{S_\var^0}
\newcommand{\varNominal}{\hat\var}
\title{\LARGE \bf Over-Approximating Minimizer Sets of Constrained Convex Programs with Parametric Uncertainty via Reachability Analysis}
\author{Brendan Gould, Chih-Yuan Chiu, Antoine P. Leeman, Kyriakos G. Vamvoudakis, Samuel Coogan, Glen Chou}
\date{Source: arXiv:2604.27355v1 (CC BY 4.0)}
\begin{document}
\maketitle

\noindent\emph{Source and licence.} \LARGE \bf Over-Approximating Minimizer Sets of Constrained Convex Programs with Parametric Uncertainty via Reachability Analysis. Version arXiv:2604.27355v1 (CC BY 4.0), distributed by arXiv under the Creative Commons Attribution 4.0 International licence, http://creativecommons.org/licenses/by/4.0/, as recorded in the arXiv licence metadata for this exact version and shown on the abstract page https://arxiv.org/abs/2604.27355v1. Full licence notice: \texttt{licenses/arxiv-2604.27355-CC-BY-4.0.txt}.\par\smallskip

\noindent\textit{Editorial scope.} Counted unit: the article's proposition Error (source0.tex lines 1555--1566). The article's complete original proof of it is delivered with it (source0.tex lines 1568--1597, one \texttt{proof} environment of 30 lines). No mathematics is written, repaired or re-derived: the statement and the proof are the article's own lines, in the article's own notation, and no retained line is substituted or re-flowed.  Retained uncounted as the source-local context the proof needs: lines 689--704; lines 719--723; lines 797--801; lines 840--853; lines 982--1002; lines 1069--1072; lines 1512--1522.  Nothing was omitted from any retained span, so the package carries no omission record.  No retained line contains a citation, so the wrapper carries no bibliography and no bibliography entry.  No retained line contains a figure environment, an image-inclusion command or drawing code, so no image is delivered and the package declares no asset dependency.  Editorial formatting only, in addition: an AMS \texttt{amsart} preamble that restates the article's own declarations and macros used by the retained lines, namely the environments \texttt{assumption}, \texttt{proposition} on separate counters, together with the standard packages those lines need.  The retained environments print in this document as Assumption 1, Proposition 1 in order of appearance, whereas the article prints the same items with the numbering of its own full document.  The complete native source of the article as distributed in the arXiv e-print for this exact version is bundled unchanged as \texttt{sources/199465/source0.tex}.
\begin{assumption}
\label{Assump: Objective Properties}
All partial second derivatives of $\obj(\var, \param)$ exist and are Lipschitz continuous. 
There exist scalars
% \antoine{scalars} 
$\strongConvexityConstant, \smoothnessConstant > 0$, with $\strongConvexityConstant \leq \smoothnessConstant$, such that $\obj(\cdot, \param)$ is $\smoothnessConstant$-smooth and $\strongConvexityConstant$-strongly convex, 
% over $\constraintSet$ 
% for each $\param \in \paramSet$, s
i.e., for each $\param \in \paramSet$ 
% \antoine{for each $\param \in \paramSet$, i.e., for any $\param \in \paramSet$; this reads weird. I think both mean the same no?}
and $\var, \var' \in \R^\varDim$:
\begin{align} \nonumber
    \obj(\var', \param) &\geq \obj(\var, \param) + \nabla \obj(\var, \param)^\top (\var' - \var) + \textstyle\frac{1}{2}\strongConvexityConstant \Vert \var' - \var \Vert_2^2, \\ \nonumber
    \obj(\var', \param) &\leq \obj(\var, \param) + \nabla \obj(\var, \param)^\top (\var' - \var) + \textstyle\frac{1}{2}\smoothnessConstant \Vert \var' - \var \Vert_2^2.
\end{align}
\end{assumption}

\begin{assumption}
\label{Assump: Parameter Set is Convex and Compact}
$\paramSet$ is convex and compact in $\R^\paramDim$. 
% Let $\diam(\paramSet) := \sup_{\param, \param'} \Vert \param - \param' \Vert_2$ denote the \textit{diameter} of $\paramSet$.
\end{assumption}

\begin{align} 
\label{Eqn: PGD Updates Notation}
    \var_{k+1} &= \dynamicsPGD(\var_k, \param, \ratePGD_k) \\ \label{Eqn: Projected Gradient Descent (PGD)}
    &:= \proj_\constraintSet\left( \var_k - \ratePGD_k \nabla_\var \obj(\var_k, \param) \right).
\end{align}

\begin{assumption}
% [\textbf{Bound on $\var_0$}]
\label{Assump: PGD Dynamics, Initial Iterate}
There is a known, bounded set $\varInitialSet \subseteq \constraintSet$ in which the initial iterate $\var_0$ of the PGD dynamics \eqref{Eqn: Projected Gradient Descent (PGD)} is fixed.
% \antoine{The verb ``drawn'' could refer to something stochastic, and could raise questions, maybe? Or I am not sure if I understand this Assumption.}
% The initial iterate $\var_0$ of the PGD dynamics \eqref{Eqn: Projected Gradient Descent (PGD)} is a fixed value drawn from a fixed, bounded set $\varInitialSet \subseteq \constraintSet$. 
% There exists
% \antoine{mismatch here: Either The initial iterate $\var_0$ is fixed, or it is an arbitrary element of a bounded set}
% \frank{Let me know if the above correction works.}
% Moreover, for some $\ratePGDMin, \ratePGDMax$ such that $0 < \ratePGDMin \leq \ratePGDMax < \frac{2}{\smoothnessConstant}$, 
% % the steplengths $(\ratePGD_k: k \geq 0)$ satisfy 
% we have 
% $\ratePGD_k \in [\ratePGDMin, \ratePGDMax]$ for each $k \geq 0$.
\end{assumption}

We consider the following dynamics over a horizon $[0, \iterationHorizon]$, with variable-parameter tuples $\state := (\var, \param)$ as states, PGD steplengths $\ratePGD$ as control inputs, and variables $\var$ as outputs:
% and 
% % initial uncertainty over the parameter
% parameter uncertainty
% over
% $\paramSet$:
% $\param_0 \in \paramSet$:
% \antoine{The $\theta_0$ is fixed no? I am not sure if it's clear if we say ``uncertainty over $\theta_0$''.}
\begin{subequations}
\label{Eqn: True Dynamics}
\begin{alignat}{2} \label{Eqn: True Dynamics, Variable}
    \var_{k+1} &= \dynamicsPGD(\var_k, \param_k, \ratePGD_k),
    \hspace{5mm} &&\forall \ k \in [0, \iterationHorizon - 1], \\ \label{Eqn: True Dynamics, Parameter}
    \param_{k+1} &= \param_k, \hspace{5mm} &&\forall \ k \in [0, \iterationHorizon - 1], \\
    \outputs_k &= \var_k, \hspace{5mm} &&\forall \ k \in [0, \iterationHorizon - 1].
\end{alignat}
\end{subequations}
% \antoine{I am not sure if we speak about $\theta^{nom}$. Is it also optimized? It could/should be. Also, it should be at least defined/mentionned.}
% \frank{$\paramNominal_0$ is mentioned immediately before \eqref{Eqn: Nominal Dynamics}. Doesn't the NLP we formulate at the end already account for the optimization over the entire nominal state-control trajectory (along with the error feedback gains through the system response), which includes $\paramNominal_0$?
% }
% \frank{In a comment, emphasize we are optimizing over nominal $\param$ as part of the NLP.}

\begin{align}
\label{Eqn: Nominal Dynamics}
    \stateNominal_{k+1} &= \dynamicsFullState(\stateNominal_k, \ratePGDNominal_k), \hspace{5mm} \forall \ k \in [0, \iterationHorizon - 1].
\end{align}

\begin{proposition}
(\textbf{$\dynamicsFullState$ Admits Lipschitz Gradients})
\label{Prop: Lipschitz Continuity of Gradient of PGD Dynamics}
Suppose Assumptions \ref{Assump: Objective Properties}, \ref{Assump: Parameter Set is Convex and Compact}, and \ref{Assump: PGD Dynamics, Initial Iterate} hold, and $\constraintSet$ is $\R^\varDim$ or an affine subspace of $\R^\varDim$.
For each $i \in [\varDim + \paramDim]$, let $\dynamicsFullState_i: \R^{\varDim + \paramDim + 1} \ra \R$ denote the $i$-th component of the dynamics $\dynamicsFullState$ in \eqref{Eqn: True Dynamics}. 
Then, $\forall \ i \in [\varDim + \paramDim]$, there exists some $\hessianBound_i > 0$ such that, $\forall \ (\state, \ratePGD), (\state', \ratePGD') \in \varConvexHull \times \paramSet \times [\ratePGDMin, \ratePGDMax]$:
\begin{align}
    \label{Eqn: Lipschitz Continuity of Gradient of PGD Dynamics}
    &\Vert \nabla \dynamicsFullState_i(\state, \ratePGD) - \nabla \dynamicsFullState_i(\state', \ratePGD') \Vert_1 \hspace{-0.4mm} \leq \hspace{-0.4mm} \hessianBound_i \Vert (\state, \ratePGD) - (\state', \ratePGD') \Vert_\infty.
\end{align}
\end{proposition}

\begin{proposition}[\textbf{Linearization Error Bound}]
\label{Prop: Quadratic Bound on Linearization Error}
\looseness-1Suppose Assumptions \ref{Assump: Objective Properties}, \ref{Assump: Parameter Set is Convex and Compact}, and \ref{Assump: PGD Dynamics, Initial Iterate} hold, and $\constraintSet$ is $\R^\varDim$ or an affine subspace of $\R^\varDim$.
For each $k \in [0, \iterationHorizon]$ and $i \in [\varDim + \paramDim]$, we define $r_{k,i}: \R^{2(\varDim + \paramDim)} \ra \R$ to be the $i$-th coordinate of the linearization error map $r_k$,
% $r_k: \R^{2(\varDim + \paramDim)} \ra \R^{\varDim + \paramDim}$,
for each $i \in [\varDim + \paramDim]$.
Then for any $\var_0 \in  \varInitialSet$, $\ratePGD_{0:k-1}, \ratePGDNominal_{0:k-1} \in [\ratePGDMin, \ratePGDMax]^k$, any $\state_k = (\var_k, \param_k) \in \varForwardReachableSet{k}(\var_0, \ratePGD_{0:k-1}) \times \paramSet$ and $\stateNominal_k = (\varNominal_k, \paramNominal_k) \in \varForwardReachableSet{k}(\var_0, \ratePGDNominal_{0:k-1}) \times \paramSet$, and any $i \in [\varDim + \paramDim]$, we have:
\begin{align}
\label{Eqn: Quadratic Bound on Linearization Error}
    |\linearizationError_{k,i}(\state_k, \ratePGD_k, \stateNominal_k, \ratePGDNominal_k)| &\leq \hessianBound_i \cdot \Vert (\stateError_k, \ratePGDError_k) \Vert_\infty^2.
\end{align}
\end{proposition}

\begin{proof}
\looseness-1 By the Mean Value Theorem, there exists a $\lambda_i \in [0, 1]$ with:
{
\setlength{\abovedisplayskip}{3.5pt}
\setlength{\belowdisplayskip}{3.5pt}
\begin{align} \nonumber
    &\dynamicsFullState_i(\state_k, \ratePGD_k) - \dynamicsFullState_i(\stateNominal_k, \ratePGDNominal_k) \\ \nonumber
    = \ &\nabla \dynamicsFullState_i\big( \lambda_i(\state_k, \ratePGD_k) + (1 - \lambda_i)(\stateNominal_k, \ratePGDNominal_k) \big)^\top (\stateError_k, \ratePGDError_k).
\end{align}
}
\hspace{-1mm}By construction, $(\state_k, \ratePGD_k), (\stateNominal_k, \ratePGDNominal_k) \in \varConvexHull \times \paramSet \times [\ratePGDMin, \ratePGDMax]$, and thus 
$\lambda_i(\state_k, \ratePGD_k) + (1 - \lambda_i)(\stateNominal_k, \ratePGDNominal_k) \in \varConvexHull \times \paramSet \times [\ratePGDMin, \ratePGDMax]$ (Note that $\varConvexHull$ and $[\ratePGDMin, \ratePGDMax]$ are convex by construction, and $\paramSet$ is convex under Assumption \ref{Assump: Parameter Set is Convex and Compact}).
% any convex combination of $(\state_k, \ratePGD_k)$ and $ (\stateNominal_k, \ratePGDNominal_k)$.
Thus, we have:
{
\setlength{\abovedisplayskip}{3.5pt}
\setlength{\belowdisplayskip}{3.5pt}
\begin{align} \nonumber
    &|\linearizationError_{k,i}(\state_k, \ratePGD_k, \stateNominal_k, \ratePGDNominal_k)| \\ \nonumber
    = \ &|\dynamicsFullState_i(\state_k, \ratePGD_k) - \dynamicsFullState_i(\stateNominal_k, \ratePGDNominal_k) - \nabla \dynamicsFullState_i(\stateNominal_k, \ratePGDNominal_k)^\top (\stateError_k, \ratePGDError_k) | \\ \nonumber
    \leq \ &\Vert \nabla \dynamicsFullState_i\big( \lambda_i(\state_k, \ratePGD_k) + (1 - \lambda_i)(\stateNominal_k, \ratePGDNominal_k) \big) - \nabla \dynamicsFullState_i(\stateNominal_k, \ratePGDNominal_k) \Vert_1 \\ \nonumber
    &\hspace{1cm} \cdot \Vert (\stateError_k, \ratePGDError_k) \Vert_\infty \\ 
    \label{Eqn: Apply Lipschitz bounds of Dynamics' Gradient}
    \leq \ &\mu_i \cdot \lambda_i \Vert (\stateError_k, \ratePGDError_k) \Vert_\infty \cdot \Vert (\stateError_k, \ratePGDError_k) \Vert_\infty,
    \\ \nonumber
    \leq \ &\hessianBound_i \cdot \Vert (\stateError_k, \ratePGDError_k) \Vert_\infty^2.
\end{align}
}
where \eqref{Eqn: Apply Lipschitz bounds of Dynamics' Gradient} follows from Prop. \ref{Prop: Lipschitz Continuity of Gradient of PGD Dynamics}.
\end{proof}

\end{document}
